The backgrounds to the $B^0 \to K^*\ell^+\ell^-$ signal
are characterised by the following categories:

\begin{itemize}
\item combinatorial background, the dominant component due
  to uncorrelated tracks and characterised by a smooth descending
  shape in the $K\pi\ell\ell$ invariant mass;
\item resonant channel feed-down: this feeds into the low-mass sideband
  and it has a significantly high branching ratio;
\item misreconstructed or partially reconstructed decays, characterised by
  real $B$ decays to final states with similar enough compositions. Most of
  these decays have missing particles in the final states and thus accumulate
  in the low-mass sideband. A detailed study of the main contributions
  from these misreconstructed $B$ decays can be found below;
\item ``peaking'' background: $B^+ \to K^+\ell^+\ell^-$ or $B^+ \to J/\psi K^+$
  will contribute to the invariant mass distribution as a peaking structure
  on the right of the signal peak. (add decays with fakes here as well?)
\end{itemize}


